\documentclass[../DoAn.tex]{subfiles}
\begin{document}

Chương này trình bày kết quả khảo sát hiện trạng các giải pháp theo dõi vị trí trẻ em hiện có trên thị trường, phân tích tổng quan chức năng hệ thống thông qua các biểu đồ use case, đặc tả chi tiết các use case quan trọng nhất, và xác định các yêu cầu phi chức năng cần đáp ứng.

\section{Khảo sát hiện trạng}
\label{section:2.1}

Các giải pháp theo dõi vị trí trẻ em hiện có thể phân thành ba nhóm chính: ứng dụng di động kết hợp điện thoại thông minh, thiết bị định vị ngắn tầm dựa trên Bluetooth/UWB, và GPS tracker chuyên dụng tích hợp kết nối di động.

\textbf{Nhóm 1 – Ứng dụng theo dõi qua điện thoại.} Life360 là đại diện tiêu biểu của nhóm này, cho phép chia sẻ vị trí giữa các thành viên trong gia đình thông qua GPS của điện thoại thông minh \cite{life360}. Ứng dụng cung cấp tính năng geofence và thông báo khi thành viên ra vào vùng an toàn. Tuy nhiên, giải pháp này hoàn toàn phụ thuộc vào việc trẻ phải mang theo điện thoại đã cài ứng dụng và có kết nối Internet, điều không khả thi với trẻ tiểu học thường bị hạn chế sử dụng thiết bị di động.

\textbf{Nhóm 2 – Thiết bị định vị ngắn tầm.} Apple AirTag và Google Find My Network sử dụng Bluetooth Low Energy (BLE) cùng công nghệ Ultra-Wideband (UWB) để xác định vị trí với độ chính xác cao trong phạm vi gần. Các thiết bị này phụ thuộc vào mạng lưới thiết bị Apple/Android lân cận để chuyển tiếp tín hiệu (crowdsourced network), khiến độ tin cậy tại khu vực thưa dân cư giảm mạnh. Ngoài ra, cơ chế này không hỗ trợ theo dõi liên tục theo thời gian thực.

\textbf{Nhóm 3 – GPS tracker chuyên dụng.} Các thiết bị như Optimus 2.0 hay Bouncie tích hợp module GPS và SIM để truyền dữ liệu qua mạng di động, cung cấp theo dõi thời gian thực mà không cần điện thoại của trẻ. Đây là hướng tiếp cận phù hợp nhất với yêu cầu đặt ra. Tuy nhiên, hầu hết các sản phẩm này đòi hỏi phí đăng ký dịch vụ hàng tháng (5–25\,USD), thiếu khả năng tùy chỉnh geofence linh hoạt, và không cung cấp API mở để tích hợp thêm.

Bảng \ref{table:comparison} tổng hợp so sánh các giải pháp hiện có theo các tiêu chí quan trọng.

\begin{table}[H]
\centering
\caption{So sánh các giải pháp theo dõi vị trí trẻ em hiện có}
\label{table:comparison}
\begin{tabular}{|l|c|c|c|c|}
\hline
\textbf{Tiêu chí} & \textbf{Life360} & \textbf{AirTag} & \textbf{GPS Tracker TM} & \textbf{Đề tài này} \\
\hline
Không cần điện thoại trẻ    & Không & Có       & Có        & Có      \\
\hline
Theo dõi thời gian thực     & Có    & Không    & Có        & Có      \\
\hline
Geofence tự động            & Có    & Không    & Hạn chế   & Có      \\
\hline
Cảnh báo đa kênh            & Có    & Không    & Hạn chế   & Có      \\
\hline
Chi phí vận hành            & Miễn phí/Thuê bao & Một lần & Thuê bao & Thấp (SIM) \\
\hline
Khả năng tùy chỉnh          & Thấp  & Rất thấp & Thấp      & Cao     \\
\hline
\end{tabular}
\end{table}

Từ phân tích trên, hệ thống cần xây dựng phải giải quyết được các điểm yếu mà các giải pháp hiện tại chưa đáp ứng đầy đủ: hoạt động độc lập không cần điện thoại của trẻ, theo dõi liên tục theo thời gian thực, geofencing linh hoạt, cảnh báo đa kênh, và chi phí vận hành thấp.

\section{Tổng quan chức năng}
\label{section:2.2}

\subsection{Biểu đồ use case tổng quát}
\label{subsection:2.2.1}

Hệ thống có ba tác nhân chính. \textbf{Phụ huynh} là người dùng chính của ứng dụng di động, thực hiện các thao tác quản lý và theo dõi. \textbf{Thiết bị IoT} (thiết bị GPS đeo trên trẻ) là tác nhân tự động, định kỳ đọc tọa độ GPS và gửi lên server mà không cần tương tác người dùng. \textbf{Dịch vụ thông báo} (Telegram Bot, Email) là tác nhân bên ngoài nhận lệnh gửi cảnh báo từ backend khi phát hiện vi phạm geofence hoặc tín hiệu SOS.

Các use case chính của hệ thống gồm: (1) Đăng ký thiết bị GPS và liên kết với tài khoản; (2) Đăng nhập hệ thống; (3) Xem vị trí trẻ theo thời gian thực trên bản đồ; (4) Tạo, chỉnh sửa và xóa vùng an toàn địa lý (geofence); (5) Nhận cảnh báo tự động khi trẻ vi phạm geofence; (6) Gửi tín hiệu SOS khẩn cấp; và (7) Xem thống kê hành trình (quãng đường, số lần vi phạm, lịch sử trên bản đồ).

\begin{figure}[H]
    \centering
    \includegraphics[width=0.85\textwidth]{Hinhve/usecase_tong_quat.png}
    \caption{Biểu đồ use case tổng quát của hệ thống}
    \label{fig:usecase_total}
\end{figure}

\subsection{Biểu đồ use case phân rã – Theo dõi và cảnh báo}
\label{subsection:2.2.2}

Use case \textit{Xem vị trí theo thời gian thực} được phân rã thành ba use case con: xem vị trí hiện tại của trẻ trên bản đồ, xem trạng thái kết nối của thiết bị (online/offline), và xem thông tin bổ sung như hướng di chuyển và tốc độ ước tính. Use case \textit{Quản lý geofence} phân rã thành: tạo vùng an toàn hình tròn bằng cách chọn tâm và bán kính trên bản đồ, đặt tên và kích hoạt hoặc vô hiệu hóa geofence, và xóa geofence không còn cần thiết.

\begin{figure}[H]
    \centering
    \includegraphics[width=0.85\textwidth]{Hinhve/usecase_phan_ra.png}
    \caption{Biểu đồ use case phân rã – Theo dõi và cảnh báo}
    \label{fig:usecase_phanra}
\end{figure}

\subsection{Quy trình nghiệp vụ}
\label{subsection:2.2.3}

Quy trình nghiệp vụ cốt lõi của hệ thống là quy trình \textbf{giám sát và cảnh báo tự động}, hoạt động như sau. Thiết bị GPS đeo trên trẻ định kỳ đọc tọa độ từ module GPS và gửi lên server qua giao thức MQTT. Backend tiếp nhận tọa độ, lưu vào cơ sở dữ liệu, sau đó kiểm tra xem tọa độ nhận được có nằm trong các vùng geofence đang hoạt động của phụ huynh hay không. Nếu trẻ đi ra ngoài vùng an toàn, hệ thống tự động gửi thông báo đến ứng dụng di động qua WebSocket đồng thời gửi cảnh báo qua Telegram Bot và Email. Song song, vị trí mới nhất được đẩy đến ứng dụng di động để cập nhật bản đồ mà không cần phụ huynh thao tác thủ công.

\section{Đặc tả chức năng}
\label{section:2.3}

\subsection{Đặc tả use case: Đăng ký thiết bị GPS}

\begin{table}[H]
\centering
\caption{Đặc tả use case Đăng ký thiết bị GPS}
\label{table:uc_device}
\begin{tabular}{|l|p{9.5cm}|}
\hline
\textbf{Tên use case}       & Đăng ký thiết bị GPS \\
\hline
\textbf{Tác nhân}           & Phụ huynh \\
\hline
\textbf{Tiền điều kiện}     & Phụ huynh đã đăng nhập; thiết bị IoT đã được lắp đặt và có mã thiết bị (device ID). \\
\hline
\textbf{Luồng sự kiện chính} &
1. Phụ huynh chọn chức năng ``Thêm thiết bị'' trong ứng dụng. \newline
2. Phụ huynh nhập mã thiết bị (device ID) được ghi trên thiết bị IoT. \newline
3. Hệ thống kiểm tra mã thiết bị chưa được đăng ký bởi tài khoản khác. \newline
4. Hệ thống liên kết thiết bị với tài khoản phụ huynh với quyền sở hữu đầy đủ. \newline
5. Ứng dụng hiển thị thiết bị mới trong danh sách và bắt đầu nhận dữ liệu vị trí từ thiết bị đó. \\
\hline
\textbf{Luồng phát sinh}    & Nếu mã thiết bị đã được đăng ký bởi tài khoản khác, hệ thống thông báo ``Thiết bị đã thuộc về tài khoản khác''. Phụ huynh có thể chia sẻ quyền theo dõi với thành viên khác trong gia đình bằng cách nhập email của họ; người được chia sẻ có thể xem vị trí nhưng không thể xóa hay chia sẻ tiếp. \\
\hline
\textbf{Hậu điều kiện}      & Thiết bị được liên kết với tài khoản phụ huynh; dữ liệu vị trí từ thiết bị bắt đầu được lưu và hiển thị trên bản đồ. \\
\hline
\end{tabular}
\end{table}

\subsection{Đặc tả use case: Đăng nhập hệ thống}

\begin{table}[H]
\centering
\caption{Đặc tả use case Đăng nhập hệ thống}
\label{table:uc_login}
\begin{tabular}{|l|p{9.5cm}|}
\hline
\textbf{Tên use case}       & Đăng nhập hệ thống \\
\hline
\textbf{Tác nhân}           & Phụ huynh \\
\hline
\textbf{Tiền điều kiện}     & Phụ huynh đã có tài khoản đăng ký trên hệ thống. \\
\hline
\textbf{Luồng sự kiện chính} &
1. Phụ huynh nhập địa chỉ email và mật khẩu vào ứng dụng. \newline
2. Hệ thống xác thực thông tin đăng nhập; nếu hợp lệ trả về token phiên làm việc. \newline
3. Ứng dụng lưu token và chuyển đến màn hình chính. \\
\hline
\textbf{Luồng phát sinh}    & Nếu thông tin không hợp lệ, hệ thống hiển thị thông báo lỗi. Nếu phiên làm việc hết hạn, ứng dụng tự động chuyển về màn hình đăng nhập. \\
\hline
\textbf{Hậu điều kiện}      & Phụ huynh đã đăng nhập thành công và có thể truy cập toàn bộ chức năng của hệ thống. \\
\hline
\end{tabular}
\end{table}

\subsection{Đặc tả use case: Xem vị trí trẻ theo thời gian thực}

\begin{table}[H]
\centering
\caption{Đặc tả use case Xem vị trí theo thời gian thực}
\label{table:uc_realtime}
\begin{tabular}{|l|p{9.5cm}|}
\hline
\textbf{Tên use case}       & Xem vị trí trẻ theo thời gian thực \\
\hline
\textbf{Tác nhân}           & Phụ huynh, Thiết bị IoT \\
\hline
\textbf{Tiền điều kiện}     & Phụ huynh đã đăng nhập; thiết bị GPS đang hoạt động và có tín hiệu GPS hợp lệ. \\
\hline
\textbf{Luồng sự kiện chính} &
1. Phụ huynh mở màn hình bản đồ trong ứng dụng. \newline
2. Ứng dụng thiết lập kết nối WebSocket đến backend và đăng ký kênh theo dõi thiết bị. \newline
3. Thiết bị IoT gửi tọa độ GPS lên topic MQTT mỗi 5 giây. \newline
4. Backend nhận tọa độ, lưu vào cơ sở dữ liệu và chuyển tiếp đến ứng dụng qua WebSocket. \newline
5. Ứng dụng cập nhật vị trí marker trên bản đồ theo mỗi bản tin nhận được. \\
\hline
\textbf{Luồng phát sinh}    & Nếu thiết bị mất tín hiệu GPS (ở trong nhà, hầm ngầm), marker hiển thị vị trí cuối cùng đã biết kèm nhãn ``Mất tín hiệu'' và thời điểm cập nhật cuối. Nếu kết nối WebSocket bị ngắt, ứng dụng tự động kết nối lại sau 3 giây. \\
\hline
\textbf{Hậu điều kiện}      & Bản đồ hiển thị vị trí mới nhất của trẻ và tiếp tục cập nhật theo thời gian thực. \\
\hline
\end{tabular}
\end{table}

\subsection{Đặc tả use case: Tạo vùng an toàn địa lý}

\begin{table}[H]
\centering
\caption{Đặc tả use case Tạo vùng an toàn địa lý}
\label{table:uc_geofence}
\begin{tabular}{|l|p{9.5cm}|}
\hline
\textbf{Tên use case}       & Tạo vùng an toàn địa lý (Geofence) \\
\hline
\textbf{Tác nhân}           & Phụ huynh \\
\hline
\textbf{Tiền điều kiện}     & Phụ huynh đã đăng nhập. \\
\hline
\textbf{Luồng sự kiện chính} &
1. Phụ huynh chọn chức năng tạo vùng an toàn trong ứng dụng. \newline
2. Phụ huynh chọn loại vùng an toàn: \newline
\hspace{1em}2a. \textit{Vùng tròn}: Phụ huynh nhấn chọn điểm trung tâm trên bản đồ và kéo thanh trượt để điều chỉnh bán kính; hệ thống cảnh báo khi trẻ vượt ra ngoài vùng tròn này. \newline
\hspace{1em}2b. \textit{Đường đi}: Phụ huynh chọn lần lượt các điểm trên bản đồ để vẽ hành trình dự kiến; hệ thống cảnh báo khi trẻ lệch khỏi hành lang cho phép xung quanh đường đi. \newline
3. Phụ huynh đặt tên cho vùng an toàn và xác nhận. \newline
4. Backend lưu cấu hình và phát sự kiện qua WebSocket; vùng an toàn hiển thị ngay trên bản đồ. \\
\hline
\textbf{Luồng phát sinh}    & Nếu phụ huynh chọn chế độ đường đi nhưng chỉ chọn được một điểm (cần tối thiểu 2 điểm để tạo đoạn đường), hệ thống không ghi nhận vùng an toàn hợp lệ; phụ huynh cần chọn thêm điểm hoặc chuyển sang chế độ vùng tròn. \\
\hline
\textbf{Hậu điều kiện}      & Geofence được lưu thành công; backend bắt đầu kiểm tra vi phạm theo cấu hình mới với mỗi điểm vị trí nhận được. \\
\hline
\end{tabular}
\end{table}

\subsection{Đặc tả use case: Nhận cảnh báo vi phạm geofence}

\begin{table}[H]
\centering
\caption{Đặc tả use case Nhận cảnh báo vi phạm geofence}
\label{table:uc_alert}
\begin{tabular}{|l|p{9.5cm}|}
\hline
\textbf{Tên use case}       & Nhận cảnh báo vi phạm geofence \\
\hline
\textbf{Tác nhân}           & Thiết bị IoT, Dịch vụ thông báo, Phụ huynh \\
\hline
\textbf{Tiền điều kiện}     & Có ít nhất một geofence đang hoạt động; thiết bị IoT đang gửi dữ liệu định kỳ. \\
\hline
\textbf{Luồng sự kiện chính} &
1. Hệ thống nhận vị trí mới từ thiết bị. \newline
2. Hệ thống kiểm tra vị trí có nằm trong vùng an toàn đang hoạt động không. \newline
3. Nếu trẻ đã ra ngoài vùng an toàn, hệ thống tạo bản ghi cảnh báo và gửi thông báo qua Telegram đến phụ huynh. \newline
4. Hệ thống đẩy thông báo cảnh báo đến ứng dụng của phụ huynh. \newline
5. Ứng dụng hiển thị cảnh báo với thông tin vị trí hiện tại và tên vùng an toàn bị vi phạm. \\
\hline
\textbf{Luồng phát sinh}    & Nếu trẻ nhanh chóng quay lại vùng an toàn, hệ thống gửi thêm thông báo ``Đã trở về vùng an toàn''. Để tránh spam, hệ thống áp dụng thời gian chờ tối thiểu 60 giây giữa hai cảnh báo liên tiếp cho cùng một geofence. \\
\hline
\textbf{Hậu điều kiện}      & Phụ huynh nhận được cảnh báo qua ít nhất một kênh thông báo. \\
\hline
\end{tabular}
\end{table}

\subsection{Đặc tả use case: Gửi tín hiệu SOS khẩn cấp}

\begin{table}[H]
\centering
\caption{Đặc tả use case Gửi tín hiệu SOS khẩn cấp}
\label{table:uc_sos}
\begin{tabular}{|l|p{9.5cm}|}
\hline
\textbf{Tên use case}       & Gửi tín hiệu SOS khẩn cấp \\
\hline
\textbf{Tác nhân}           & Thiết bị IoT, Dịch vụ thông báo, Phụ huynh \\
\hline
\textbf{Tiền điều kiện}     & Thiết bị IoT đang hoạt động và có kết nối MQTT đến server. \\
\hline
\textbf{Luồng sự kiện chính} &
1. Trẻ em nhấn nút SOS trên thiết bị. \newline
2. Thiết bị đọc tọa độ GPS hiện tại và gửi tín hiệu SOS kèm vị trí lên hệ thống. \newline
3. Hệ thống nhận tín hiệu SOS và tạo bản ghi cảnh báo khẩn. \newline
4. Hệ thống gửi ngay thông báo SOS qua Telegram đến phụ huynh kèm vị trí và đường link bản đồ. \newline
5. Hệ thống đẩy cảnh báo đến ứng dụng của phụ huynh. \newline
6. Ứng dụng hiển thị popup cảnh báo khẩn và vị trí hiện tại trên bản đồ. \\
\hline
\textbf{Luồng phát sinh}    & Nếu thiết bị chưa có tín hiệu GPS tại thời điểm nhấn nút, bản tin SOS vẫn được gửi lên kèm tọa độ cuối cùng đã biết và ghi chú ``Vị trí có thể không chính xác''. Nếu kết nối MQTT bị mất, thiết bị lưu sự kiện SOS vào bộ nhớ và gửi lại ngay khi kết nối được khôi phục. \\
\hline
\textbf{Hậu điều kiện}      & Phụ huynh nhận được cảnh báo SOS qua tất cả các kênh thông báo đã cấu hình; sự kiện SOS được lưu vào lịch sử để tra cứu sau. \\
\hline
\end{tabular}
\end{table}

\subsection{Đặc tả use case: Xem thống kê hành trình}

\begin{table}[H]
\centering
\caption{Đặc tả use case Xem thống kê hành trình}
\label{table:uc_stats}
\begin{tabular}{|l|p{9.5cm}|}
\hline
\textbf{Tên use case}       & Xem thống kê hành trình \\
\hline
\textbf{Tác nhân}           & Phụ huynh \\
\hline
\textbf{Tiền điều kiện}     & Phụ huynh đã đăng nhập; hệ thống đã có dữ liệu vị trí và cảnh báo được lưu trữ. \\
\hline
\textbf{Luồng sự kiện chính} &
1. Phụ huynh chọn chức năng thống kê và chọn khoảng thời gian cần xem. \newline
2. Hệ thống truy vấn dữ liệu vị trí và cảnh báo trong khoảng thời gian đã chọn. \newline
3. Ứng dụng hiển thị các chỉ số tổng hợp: tổng quãng đường di chuyển và số lần ra khỏi vùng an toàn. \newline
4. Ứng dụng vẽ hành trình trên bản đồ thể hiện toàn bộ đường đi của trẻ trong khoảng thời gian đó. \\
\hline
\textbf{Luồng phát sinh}    & Nếu không có dữ liệu cho khoảng thời gian đã chọn, ứng dụng hiển thị các chỉ số bằng 0 và bản đồ trống. \\
\hline
\textbf{Hậu điều kiện}      & Phụ huynh nắm được tổng quãng đường trẻ đã di chuyển, số lần vi phạm vùng an toàn, và hành trình cụ thể trên bản đồ. \\
\hline
\end{tabular}
\end{table}

\section{Yêu cầu phi chức năng}
\label{section:2.4}

\textbf{Hiệu năng.} Độ trễ end-to-end từ khi thiết bị gửi tọa độ đến khi ứng dụng cập nhật trên bản đồ không vượt quá 5 giây trong điều kiện mạng di động thông thường (3G/4G). Backend phải có khả năng xử lý đồng thời ít nhất 100 kết nối WebSocket và xử lý tối thiểu 50 bản tin MQTT mỗi giây.

\textbf{Độ tin cậy.} Thiết bị phải tự động kết nối lại MQTT sau khi mất kết nối mà không cần can thiệp thủ công, với thời gian phục hồi không quá 30 giây. Dữ liệu vị trí được gửi định kỳ qua MQTT; khi kết nối bị gián đoạn, thiết bị tự động kết nối lại và tiếp tục gửi, nên việc thất lạc một vài bản tin lẻ trong thời gian gián đoạn không ảnh hưởng đến khả năng theo dõi liên tục.

\textbf{Tính dễ sử dụng.} Giao diện ứng dụng di động phải đủ đơn giản để phụ huynh ít kinh nghiệm công nghệ có thể sử dụng mà không cần hướng dẫn. Tất cả thao tác quan trọng phải thực hiện được trong không quá 3 lần nhấn từ màn hình chính.

\textbf{Bảo mật.} API backend xác thực người dùng bằng cơ chế JSON Web Token (JWT): sau khi đăng nhập thành công, server phát hành token JWT có ký số và thời hạn, ứng dụng đính kèm token này vào header \texttt{Authorization} của mỗi yêu cầu. Mật khẩu được lưu trữ dưới dạng băm bcrypt, không lưu plaintext. Cơ sở dữ liệu MongoDB áp dụng xác thực tài khoản để giới hạn truy cập.

\textbf{Yêu cầu kỹ thuật.} Backend triển khai trên môi trường đám mây với uptime tối thiểu 99\%. Ứng dụng di động hỗ trợ Android phiên bản 6.0 (API 23) trở lên và iOS phiên bản 12.0 trở lên. Dữ liệu vị trí được lưu trữ tối thiểu 30 ngày trước khi xóa tự động.

Kết chương: Chương 2 đã phân tích chi tiết hiện trạng các giải pháp theo dõi vị trí trẻ em, chỉ ra các hạn chế còn tồn tại và xác định rõ yêu cầu cho hệ thống cần xây dựng. Kết quả khảo sát cho thấy hướng tiếp cận thiết bị IoT chuyên dụng kết hợp GPS và kết nối di động là phù hợp nhất. Các yêu cầu chức năng và phi chức năng được đặc tả trong chương này sẽ là cơ sở để lựa chọn công nghệ và thiết kế kiến trúc hệ thống, được trình bày trong các chương tiếp theo.

\end{document}
